\section{Synthetic recovery experiment}
\label{sec:recovery}

% Loads the recovery macros (\thetaRecoveryCorr, \difficultyMinKnots, ...) and the
% results table -- both generated from stats.json. Input early so the macros are
% defined before the prose cites them; the table floats to the top of a page.
% !!! GENERATED FILE -- DO NOT EDIT BY HAND !!!
% Produced by reproducibility/gen_results_tex.py from the stats.json that
% `inverse-suitability latent-demo` writes. Regenerate with `make repro`.
% Every value below is a RECOVERED estimate; ground-truth DGP parameters are
% stated as fixed prose in the surrounding section, never here.
%
% Inline macros (used in 05-synthetic-recovery.tex prose):
\newcommand{\thetaRecoveryCorr}{0.96}
\newcommand{\thetaRecoveryCorrFull}{0.9638}
\newcommand{\difficultyMinKnots}{16.7}
\newcommand{\discriminationA}{2.16}
\newcommand{\heldOutBSS}{0.329}
% signed, WITHOUT math delimiters -- wrap in $...$ at the call site.
\newcommand{\heldOutBSSsigned}{+0.33}
\newcommand{\nTrain}{2,400}
%
\begin{table}[t]
  \centering
  \caption{Synthetic recovery on the reference cohort ($80$ riders $\times\,30$
  paired outcomes, $n_{\text{train}} = 2,400$). Every value is recovered
  from behavioural outcomes alone by marginal maximum likelihood; none is supplied
  to the fit. The data-generating difficulty is minimised at the ground-truth
  optimum $x^\star = 18$~kn with true discrimination $a_{\text{true}} = 2.0$
  (stated in the text, not fit here). Reproduced by a single command
  (\texttt{inverse-suitability latent-demo}); see Appendix.}
  \label{tab:recovery}
  \begin{tabular}{@{}llr@{}}
    \toprule
    Quantity & Symbol & Recovered \\
    \midrule
    Latent-skill recovery (Pearson) & $\mathrm{corr}(\theta,\hat\theta)$ & 0.964 \\
    Difficulty-minimum location     & $\hat{x}^\star$ (kn)               & 16.7 \\
    Discrimination                  & $\hat{a}$                          & 2.16 \\
    Held-out Brier skill score      & $\mathrm{BSS}$ vs.\ single curve   & $+0.329$ \\
    Marginal matches single curve   & --                                    & \checkmark \\
    Identification gate passed      & --                                    & \checkmark \\
    \bottomrule
  \end{tabular}
\end{table}


Because no real cohort is used, the validity of the method rests on a
\emph{synthetic recovery} study: we generate data from a known ground truth, fit
the model, and measure how well it recovers that truth. This is the standard
validation for a methods paper, and it is the same role a simulation study plays
throughout statistics---it isolates the estimator's behaviour from the confounds
of any particular dataset.

\subsection{Data-generating process}
\label{sec:dgp}

The generator draws a two-factor cohort with \emph{known} skill and difficulty.
Ground-truth skills are $\skill_\rider \sim \mathcal{N}(0,1)$, one per rider. The
ground-truth difficulty is a parabola in difficulty space,
\begin{equation}
  \label{eq:dgp-delta}
  \diff(\metric) = \left(\frac{\metric - \metric^\star}{w}\right)^{2} + f,
  \qquad \metric^\star = 18~\text{kn},\; w = 8,\; f = -1,
\end{equation}
so that difficulty is \emph{minimised at the ground-truth optimum
$\metric^\star = 18$~knots} and rises symmetrically toward calm and storm---a
single-peaked suitability profile once passed through the link. Each rider is
observed across conditions drawn uniformly over $[3, 35]$ knots (guaranteeing a
connected incidence graph, Section~\ref{sec:graph}), and outcomes are Bernoulli
draws from \eqref{eq:model} with \emph{true discrimination}
$a_{\text{true}} = 2.0$. The reference cohort has $80$ riders $\times\,30$
outcomes each. These parameters---$\metric^\star = 18$~kn, $a_{\text{true}} = 2.0$,
and the cohort shape---are the ground truth supplied to the \emph{generator};
they are never supplied to the \emph{fit}. The recovered quantities reported below
are what the estimator infers from the outcomes alone.

\subsection{Metrics and results}

We assess recovery on three axes: (i) latent-skill recovery, the Pearson
correlation between the true $\skill_\rider$ and the estimated $\hat\skill_\rider$;
(ii) difficulty-shape recovery, summarised by the estimated location of the
difficulty minimum $\hat\metric^\star = \arg\min_\metric \hat\diff(\metric)$,
which should land near the true $18$~kn; and (iii) predictive gain, the held-out
Brier Skill Score of the skill-conditioned model against the single-curve
baseline under a stratified $80/20$ split. We also verify the backward-compatible
identity of Section~\ref{sec:marginal}---that integrating skill out of the fitted
model matches a single-curve fit---and that the identification gate of
Section~\ref{sec:graph} passes on this (deliberately connected) cohort.

Table~\ref{tab:recovery} reports the outcome. The model recovers latent skill at
correlation $\thetaRecoveryCorr$ against ground truth; it locates the difficulty
minimum at $\hat\metric^\star = \difficultyMinKnots$~knots, within $\pm 3$ knots of
the true $18$-knot optimum; and it improves the held-out Brier Skill Score by
$\heldOutBSSsigned$ over the single-curve baseline, confirming that
skill-conditioning earns its added complexity out of sample rather than merely
fitting the training data more flexibly. The recovered discrimination is
$\hat{a} = \discriminationA$ (true $a_{\text{true}} = 2.0$), the marginal matches
the single-curve fit, and the identification gate passes, all on
$n_{\text{train}} = \nTrain$ training outcomes.

\subsection{Reproducibility}
\label{sec:repro}

The experiment is a single command. Every value in Table~\ref{tab:recovery} is
emitted by
\begin{quote}
  \verb|uv run inverse-suitability latent-demo --out DIR|
\end{quote}
which fits the reference cohort and writes a \texttt{stats.json}; the paper's
table and inline figures are rendered directly from that file, so the reported
numbers cannot drift from the code (Appendix~\ref{app:repro}). The run is
deterministic (fixed cohort and fit seeds), completes in well under a minute on a
laptop, touches no proprietary data, and requires no network access. The
data-generating process is fully specified by \eqref{eq:dgp-delta} and
Section~\ref{sec:dgp}; a reviewer can regenerate the cohort and re-fit
independently.
